\documentclass{article}
\usepackage{iclr2026_conference,times}

\input{math_commands.tex}

\usepackage{hyperref}
\usepackage{url}

\title{Dexterous Piano Track Proposal: Improving PianoMime with Various Tricks}

\author{Sining Zhoubian \\
ID: 2025210759 \\
\texttt{zbsn21@mails.tsinghua.edu.cn} \\
\And Kun Feng \\
ID: 2025316279 \\
\texttt{fk25@mails.tsinghua.edu.cn} \\
\And Xingyu Shen \\
ID: 2025312187 \\
\texttt{xingyu-c21@mails.tsinghua.edu.cn} \\
\And Wenkai Li \\
ID: 2025316137 \\
\texttt{lwk21@mails.tsinghua.edu.cn} \\
\And Yoshua Bengio \\
ID: 0 \\
\texttt{bengio21@mails.tsinghua.edu.cn} \\
}

\newcommand{\policy}{\pi}
\newcommand{\residual}{\Delta a}

\iclrfinalcopy
\begin{document}

\maketitle

\begin{abstract}
Piano playing with dexterous robot hands is a useful benchmark for long-horizon visuomotor imitation, precise contact timing, and high-dimensional control. This proposal studies how to improve PianoMime, a framework that learns a generalist piano-playing policy from internet demonstrations and then adapts to individual songs with residual reinforcement learning. The current pipeline is capable, but still can not perform well on some songs. To strengthen the RL pipeline, we propose several reinforcement tricks from pretraining and reward shaping to goal refinement. The expected outcome is a more stable optimization path that preserves PianoMime's broad imitation prior while improving key-press accuracy, temporal alignment, and robustness on held-out songs.
\end{abstract}

\section{Related Work}

PianoMime aims to control two dexterous Shadow Hands in RoboPianist so that they can play piano pieces from MIDI goals and human demonstration videos \citep{qian2024pianomime,zakka2023robopianist}. Its key idea is to combine broad offline imitation with narrow online correction. The offline part learns from many internet demonstrations, while the online part adapts to a specific song in simulation.

The baseline can be summarized as a two-stage control stack. First, PianoMime trains diffusion policies to model demonstration actions. In the provided implementation, the high-level model maps observations and MIDI-conditioned context to fingertip trajectories, and the low-level model maps richer state and trajectory context to hand control actions. Both models use a conditional 1D U-Net trained as a denoising diffusion model: noisy action sequences are sampled at random diffusion timesteps, and the network predicts the injected noise with an MSE objective. This gives a multi-task generative prior over plausible piano-playing motions.

Second, PianoMime trains song-specific residual policies with PPO. The residual environment loads left- and right-hand demonstration trajectories, computes an inverse-kinematics prior action from the current demonstration frame, and adds the learned policy output as a residual before stepping the simulator. In simplified form, the executed hand action is
\begin{equation}
    a_t = a^{\mathrm{prior}}_t + \alpha \, \residual_t,
\end{equation}
where $a^{\mathrm{prior}}_t$ is obtained from demonstration-guided IK, $\residual_t \sim \policy_{\theta}(\cdot \mid o_t)$ is the PPO policy output, and $\alpha$ is the residual scale. This design is sample-efficient because PPO does not need to discover piano-playing behavior from scratch; it only needs to compensate for tracking, contact, and timing errors.

However, the code-level connection between the two stages is indirect. The residual policy is not necessarily trained on actions generated by the diffusion model; it is typically trained with stored demonstration trajectories and the residual wrapper. Therefore, the Stage-1 diffusion policy and Stage-2 residual policy can be viewed as two separately optimized modules: one learns a broad imitation prior, and the other learns closed-loop correction around a prior trajectory. This separation motivates the proposed improvements below.

\section{Proposed Ideas for Optimization}

\subsection{Residual-Aware Imitation}

The first idea is to align Stage 1 and Stage 2 by making residual learning part of the imitation interface. Instead of treating the diffusion policy output and the PPO residual as independent quantities, we propose to train or fine-tune a policy around the composed action
\begin{equation}
    a_t = a^{\mathrm{diff}}_t + \alpha \, \policy_{\theta}(o_t, a^{\mathrm{diff}}_t, g_t),
\end{equation}
where $a^{\mathrm{diff}}_t$ is the action or target trajectory generated by the diffusion model, and $g_t$ contains note, fingering, and lookahead information. Practically, this can start with a minimal change: include the prior action and the short-horizon demonstration or diffusion trajectory as explicit observations for PPO, then evaluate whether the residual policy trained around the diffusion prior transfers better than one trained only around stored demonstrations.

The justification is that the residual policy should learn the error distribution of the prior that will actually be used at inference time. If PPO is trained around human-extracted trajectories but evaluated around generated diffusion trajectories, the residual may overfit to a cleaner or differently distributed prior. Residual-aware imitation reduces this train--test mismatch and gives a clear optimization target: Stage 1 proposes feasible motion, while Stage 2 learns closed-loop corrections conditioned on that proposal.

\subsection{Reward Shaping for Timing, Contact, and Smoothness}

The second idea is to refine the reward beyond raw note success. Piano playing failures are often not binary: the robot may hit the correct key too early, press with the wrong finger, miss sustain timing, or reach the target with excessive jitter. We are considering to incorporate a shaped reward with four additional terms:
\begin{equation}
    r_t = r^{\mathrm{base}}_t + \lambda_1 r^{\mathrm{timing}}_t + \lambda_2 r^{\mathrm{pre}}_t - \lambda_3 c^{\mathrm{smooth}}_t - \lambda_4 c^{\mathrm{collision}}_t.
\end{equation}
The timing term rewards key activation near the MIDI onset window; the pre-positioning term rewards fingertips that approach upcoming assigned keys before onset; the smoothness cost penalizes large action changes or high fingertip acceleration; and the collision cost discourages hand--hand or finger--finger interference when collision information is available.
This is expected to improve learning because the existing sparse success signal gives delayed feedback for errors that originate many frames earlier. For example, a missed note may be caused by poor approach direction, late pre-positioning, or unstable residual actions. Dense shaping gives PPO more local gradients while preserving the final objective.

\subsection{Contact-Centric Goal Representation}

The third idea is to redesign the goal representation from symbolic note prediction into future contact planning. Instead of conditioning the policy only on current MIDI states or future note embeddings, we propose to represent piano playing as a sequence of future contact events.

More specifically, the future goal is represented as
\begin{equation}g_t={c_1, c_2, \dots, c_N},\end{equation}
where each contact token is defined as
\begin{equation}c_i=(k_i,f_i,\tau_i,\rho_i).
\end{equation}
Here, $k_i$ denotes the target piano key, $f_i$ denotes the assigned finger, $\tau_i$ denotes the desired contact time, and $\rho_i$ denotes a learned or heuristic urgency score.

The urgency term is designed to prioritize imminent contacts:
\begin{equation}
\rho_i=\exp(-\gamma (\tau_i - t)),
\end{equation}
where $\gamma$ controls temporal decay. Contacts with smaller remaining time therefore receive higher priority.

The residual policy is conditioned on the contact representation through an attention mechanism:
\begin{equation}a_t=a_t^{\mathrm{prior}}+\alpha\pi_\theta(o_t,\mathrm{Attn}(o_t, g_t)).\end{equation}

This formulation improves learning by explicitly modeling future interaction constraints rather than only symbolic note sequences. Piano playing is fundamentally a contact scheduling problem involving precise timing, finger occupancy, and anticipatory motion preparation. By representing future contacts directly, the policy can better reason about finger transitions, chord preparation, and temporal conflicts before note onset.

Compared with raw MIDI conditioning, the proposed representation provides a more physically grounded objective and reduces ambiguity in long-horizon dexterous control.

\subsection{RL Tricks for Single-Song Policies}

Fine-tuning residual policies on specific, long piano pieces challenges standard PPO due to compounding errors and sparse exploration. To stabilize and accelerate this single-song optimization, we propose four targeted reinforcement learning tricks. First, Reference State Initialization (RSI) randomly initializes episodes at different MIDI keyframes to prevent early-stage overfitting and ensure uniform exploration across the entire song. Second, a Tempo-based Curriculum scales down the playback speed during early training (e.g., $0.5\times$), allowing the policy to master spatial key-press accuracy before adapting to strict temporal constraints. Third, an Asymmetric Actor-Critic architecture feeds privileged simulation states (e.g., exact contact forces, infinite MIDI lookahead) solely to the critic, which reduces value estimation variance without violating inference-time observation constraints. Finally, Residual Action Regularization applies a magnitude penalty to the PPO output $\residual_t$. This forces the policy to make minimalist corrections, preserving the smooth and human-like kinematic prior generated by the diffusion model.

\subsection{Music Language Model Perplexity as an Auxiliary Reward}

The fifth idea introduces a learned, data-driven critic of musical quality and uses its likelihood as an auxiliary reward for the single-song residual policy. The existing PPO rewards only check whether each MIDI event is hit at the right time and finger; they cannot tell whether the resulting sequence is musically coherent, so policies producing small timing offsets, ghost keypresses, or local note swaps can still satisfy them while degrading the perceived performance.

We pretrain a music language model (MusicLM) on a classical piano MIDI corpus by tokenizing each piece into (pitch, duration, velocity, time-shift) events and training an autoregressive Transformer with a next-token objective. The frozen model defines a perplexity
\begin{equation}
    \mathrm{PPL}_\phi(x_{1:T}) = \exp\!\bigg(-\frac{1}{T}\sum_{t=1}^{T}\log p_\phi(x_t \mid x_{<t})\bigg)
\end{equation}
on any MIDI sequence, where lower values indicate that the sequence is more typical of classical music. During single-song PPO training, we tokenize the simulator's MIDI events on a sliding window of $W$ steps and add a bounded auxiliary reward
\begin{equation}
    r^{\mathrm{ppl}}_t = -\beta\bigl(\log \mathrm{PPL}_\phi(x_{t-W:t}) - \ell_0\bigr)
\end{equation}
to the shaped reward proposed above with weight $\lambda_5$, where $\ell_0$ is the log-perplexity of the target song's ground-truth MIDI and $\beta$ scales the term. Because the MusicLM is frozen, this leaves the IK-prior plus residual interface $a_t = a^{\mathrm{prior}}_t + \alpha\,\residual_t$ unchanged, while penalising unmusical artefacts that per-step rewards miss and supplying long-horizon, sequence-level supervision. The same $\mathrm{PPL}_\phi$ also serves as a song-agnostic evaluation metric alongside note F1.

\section{Expected Contributions}

This proposal is intentionally scoped as an optimization study rather than a new full system. The expected contributions are: (1) a clear analysis of the PianoMime pipeline, and relevant baseline performance results; (2) an enhanced robopianist pipeline based on PianoMime including some optimization techniques and their actual performance; and (3) insights on improving a robopianist pipeline.

\bibliography{iclr2026_conference}
\bibliographystyle{iclr2026_conference}

\end{document}
